% B05 / A023 multifactor point-defect refund theorem
\section*{Ordered Möbius ownership for disjoint point defects}

Let $q_n$ be $+1$ at one point of the uniform $n$-cube and $-1$
elsewhere.  For a finite index set $A$ take the factors $q_{n_i}$ on
pairwise disjoint physical coordinates and put $x_i=2^{-n_i}$.  For
$B\subseteq A$, define
\[
 p_B=\frac{1-\prod_{i\in B}(1-2x_i)}2,
 \qquad R(B)=F(p_B),\qquad R(\varnothing)=0,
\]
where the complete selector packet against a constant is
\[
 F(p)=4p(1-p)-2(1-p)^2\log_2(1-p)-2p^2\log_2p
 +(1-2p)^2\log_2|1-2p|.
\]
The zero-log convention supplies all endpoints.  The scalar formula is the
independently reviewed B19 reconstruction of the complete conditional-label
plus phase packet.

For nonempty $B$ define its Möbius interaction
\[
 \Omega(B)=\sum_{C\subseteq B}(-1)^{|B|-|C|}R(C).
\]
Then Möbius inversion gives
\begin{equation}
 R(A)=\sum_{\varnothing\ne B\subseteq A}\Omega(B).       \tag{1}
\end{equation}

Fix any deterministic factor order $\pi=(i_1,\ldots,i_r)$ and let
$A_j=\{i_j,\ldots,i_r\}$.  Assign an interaction $\Omega(B)$ to its earliest
member in this order.  The total assigned to $i_j$ is exactly
\begin{equation}
 \sum_{\substack{B\subseteq A_j\\ i_j\in B}}\Omega(B)
 =R(A_j)-R(A_{j+1}).                                    \tag{2}
\end{equation}
Indeed, subtract the two Möbius inversions for $A_j$ and $A_{j+1}$.
Thus every interaction of every order is owned exactly once.

This ownership is realized by the physical factor-first tree.  If the
current first factor has residual dimension $m$ and the unrevealed remaining
factors are $A_{j+1}$, define
\[
 \Gamma_{m\mid A_{j+1}}
 =R(\{q_m\}\cup A_{j+1})
 -\frac12R(\{q_{m-1}\}\cup A_{j+1})
 -\frac12R(A_{j+1}),                                    \tag{3}
\]
where $q_0$ is deleted, so the two terms at $m=1$ both equal
$R(A_{j+1})$.  The continuing child at level $m$ is reached with probability
$2^{m-n_{i_j}}$.  Direct cancellation gives
\begin{equation}
 \sum_{m=1}^{n_{i_j}}2^{m-n_{i_j}}\Gamma_{m\mid A_{j+1}}
 =R(A_j)-R(A_{j+1}).                                    \tag{4}
\end{equation}
Equations (2)--(4) identify the physical-tree phase with earliest-factor
Möbius ownership.

Let $P_j$ and $N_j$ be the reach-weighted positive and negative parts of the
credits in (3).  Then
\[
 P_j-N_j=R(A_j)-R(A_{j+1}).
\]
Summing over all factors and using $R(A_{r+1})=0$ yields the complete
single-use identity
\begin{equation}
 \boxed{R(A)+\sum_{j=1}^rN_j=\sum_{j=1}^rP_j.}           \tag{5}
\end{equation}
No ancestor credit, terminal packet or higher-order interaction is reused.
Ownership depends on the chosen order; no canonical symmetric allocation is
claimed.

\subsection*{Higher interactions are not product-small}

For $r$ equal factors of dimension $n$, write $\delta=2^{-n}$ and
\[
 \Omega_r(n)=\sum_{j=1}^r(-1)^{r-j}\binom rj
 F\!\left(\frac{1-(1-2\delta)^j}{2}\right).
\]
The exact scalar expansion is
\[
 F(p)=4p+p^2\left(2\log_2\frac1p-4+\frac3{\ln2}\right)
 +O\!\left(p^3\log\frac1p\right).
\]
For every fixed $r\ge3$, polynomial finite differences cancel the linear
term through order $\delta^2$, the $\log(1/\delta)$ term, and the remaining
quadratic polynomials.  Therefore
\begin{equation}
 \boxed{
 \Omega_r(n)=c_r\delta^2+O_r\!\left(\delta^3\log\frac1\delta\right),
 \qquad
 c_r=-2\sum_{j=1}^r(-1)^{r-j}\binom rj j^2\log_2j.}       \tag{6}
\end{equation}
In particular
\[
 c_3=24-18\log_2 3=-4.529325\ldots<0.
\]

To determine the sign for all $r\ge3$, put $\phi(x)=x^2\ln x$ and
$\phi(0)=0$.  Its $r$th forward difference at zero is the sum in (6),
up to $1/\ln2$.  The repeated-integral representation of a forward
difference and
\[
 \phi^{(r)}(x)=2(-1)^{r-3}(r-3)!x^{2-r}
\]
show that this difference has sign $(-1)^{r-3}$; the singularity is
integrable over the $r$-cube.  Consequently
\begin{equation}
 \operatorname{sgn}(c_r)=(-1)^r.                         \tag{7}
\end{equation}

Thus every fixed interaction order survives at scale $\delta^2$, with
alternating sign.  A pair-only ledger, or a rule pricing an $r$-fold
interaction by $O(\delta^r)$, is false.  Equation (5) nevertheless assigns
all such interactions once.  What remains open is a dimension-independent
bound on the total positive credits (equivalently the parent packet plus all
negative credits) by nonreused influence, uniformly over the number and sizes
of factors, and then beyond disjoint point defects to arbitrary branching and
shared-variable trees.
